\scanpage{30}\mspage{28}
3. \uline{Passage aux catégories de diagrammes}

Nous \deleted{allons} convenons de noter la catégorie $X$
$\unsure{M|\mathcal{M}}$, \unsure{car} par la suite elle jouera le rôle d'une
``catégorie de modèles''. On désignera par contre
par $X$ les ``catégories d'indices'' sur lesquelles
on veut construire les \deleted{\unsure{fib. de}} ``catégories de
diagrammes'' de type $X$
\[
M(X)=\underline{\mathrm{Hom}}(X,M)\quad .
\]
On se donne donc
\diagram{\begin{tikzcd}[ampersand replacement=\&,row sep=1ex,column sep=0.4em]
\unsure{\Sigma|Z}\arrow[d] \& {}\subset\mathrm{Fl}(M), \& \Sigma\supset\text{ens. des flèches identiques},\\
\text{\uline{quasi-isomorphismes}}\&{}\&{}
\end{tikzcd}}
et pour \unsure{tte} catégorie $X$, on définit
\[
\Sigma_{\inserted{X}}\deleted{(\unsure{X})}\subset\mathrm{Fl}\,M(X)
\]
\unsure{ainsi} : Si $F,G\in\mathrm{Ob}\,M(X)$ ($F,G\colon X\rightrightarrows M$),
une flèche $u\colon F\unsure{\mathord{-}|\to}G$ (\unsure{hom|Hom} de foncteurs)
est dans $\Sigma_X\deleted{(\unsure{X})}$ ssi $\forall\,x\in X$, $u(x)\colon F(x)\mathord{-}G(x)$
est dans $\Sigma$ (``critère \inserted{de quasi-isomorphie} fibre par fibre'', pour
\inserted{l'hom. entre} les copréfaisceaux $F,\unsure{G|K}$ envisagés).

Soient
\[
a,b\in\mathrm{Ob}\,M
\]
et considérons les foncteurs constants correspondants
sur $X$
\[
a_X,b_X\in\mathrm{Ob}\,M(X),
\]
% archive sheet number: 29
